\section*{Limitations}
\label{sec:limitations}

\paragraph{Scale and variance.} The complete results are on one scheme and language, Chinese UMR, with \nDocs{} documents and \nSentences{} sentences in the sequential run; three documents have only 10--12 sentences. Each configuration was run once, in one document order; we report paired tests over documents and a sentence-level bootstrap, but not run-to-run variance of building or annotating \gap{G-var}.

\paragraph{No test before keeping an edit.} \sys keeps every edit that applies cleanly; its quality is measured only on the next document. Methods that test edits on held-in data may avoid some of the failure modes we report, and we have not compared them \gap{G9}.

\paragraph{Human influence.} The seed harness, the judge's checks and the builder's prompts were written by the authors, and one rule step was changed by hand during the run (\S\ref{sec:harness}). The suggestions in \S\ref{sec:analysis} were coded by one person.

\paragraph{Models.} The main run uses one closed model family for both annotating and building. Results with other models (Table~\ref{tab:models}) use a different harness and fewer documents, and DeepSeek ran without thinking.

\paragraph{Data.} The leakage guard blocks verbatim copying of eight characters or more, not paraphrase or document-specific knowledge. The source texts of the documents may appear in pretraining data; \gap{state whether the Chinese UMR gold is public}. Gold conventions differ between documents (\S\ref{sec:analysis}), which bounds what any learner can reach.

\section*{Ethics Statement}
Automating annotation can reduce the paid work available to annotators. We see the learned harness as a way to produce drafts and to audit gold, with people deciding conventions; its suggestions about inconsistent gold are addressed to them. All model calls were made through commercial APIs or open-weight models under their terms of use; the UMR documents are used for research under their release terms.
